% Transcription — fonds Alexandre Grothendieck, shelfmark 158, batch 2 (pages 21-40).
% Established from the facsimile archives/batches/158/batch-02.pdf (a mirror of
% https://grothendieck.umontpellier.fr/158.pdf), per the skill
% transcribe-grothendieck. No cover sheet inside the batch file: PDF page N =
% archivists' page N+20, checked against the pencilled numbers bottom-left on
% the first, second and last leaves, which read 21, 22 and 40.
% His own pagination runs 3 to 18 across these leaves: archivists' page N carries
% his page N-18 — except that leaf 35 carries "16" and leaf 36 carries "15 bis",
% so those two are filed in the reverse of his order (15 bis breaks off exactly
% where 16 resumes). The run is transcribed below in the archivists' order.
% Pages 34, 38 and 40 are unrelated versos (a blank Inspection Académique
% letterhead, then twice a circular from the same office) and are absent.
% That circular is dated Avignon, 6 December 1990, which agrees with the
% inventory's [vers 1990-1991]; it is not recorded in the edition itself.
% Pass: Opus 5 (claude-opus-5), 2026-09-20 — first pass, unchecked against the pages by a human.
\documentclass[11pt,a4paper]{article}
% Le préambule commun à toutes les transcriptions du fonds Grothendieck.
%
% Il tient en une idée : **l'incertitude se compose, elle ne se résout pas.**
% Une transcription qui lisse un mot douteux a détruit la seule chose pour
% laquelle on la faisait — un lecteur doit pouvoir savoir, à chaque endroit,
% ce qui était sur le feuillet et ce qui a été ajouté. Les six macros qui
% suivent sont donc l'appareil critique, et non de la décoration.
%
% Les mêmes commandes sont reconnues par `scripts/render.mjs`, qui produit la
% vue de lecture du site. Une macro ajoutée ici doit l'être aussi là-bas,
% faute de quoi le rendu échouera bruyamment — ce qui est voulu : un rendu
% silencieusement faux est pire qu'un rendu absent.

\ProvidesPackage{grothendieck}[2026/08/08 Transcriptions du fonds Grothendieck]

\usepackage{amsmath,amssymb,amsthm}
\usepackage{mathrsfs}
\usepackage{stmaryrd}
% Les diagrammes commutatifs sont partout dans ce fonds : c'est la langue
% même de la géométrie algébrique de Grothendieck, et une transcription qui
% les rendrait en prose ne transcrirait rien.
\usepackage{tikz-cd}
% Français par défaut — c'est la langue des feuillets — mais l'anglais est
% chargé aussi : la traduction et le résumé sont des documents anglais, et la
% typographie française (espaces avant les deux-points) y serait une faute.
% Ces documents basculent par \selectlanguage{english} en tête de corps.
\usepackage[english,french]{babel}
\usepackage{fontspec}
\usepackage{geometry}
\geometry{a4paper,margin=25mm,marginparwidth=28mm}
\usepackage{marginnote}
\usepackage{xcolor}
% ulem, not soul. `soul`'s \st and \ul are built on a character-by-character
% scan that XeLaTeX's font machinery breaks: the document fails with an
% undefined control sequence rather than degrading. ulem does the same two jobs
% and is Unicode-engine safe. `normalem` stops it hijacking \emph, which the
% transcriptions use for Grothendieck's own emphasis.
\usepackage[normalem]{ulem}
\usepackage{hyperref}
\hypersetup{colorlinks=true,linkcolor=black,urlcolor=black,pdfborder={0 0 0}}

% --- Métadonnées du lot -----------------------------------------------------
% Reprises telles quelles par le script de rendu, qui les affiche en tête de
% la vue de lecture. Elles fixent ce que le fichier prétend couvrir : sans
% elles, une transcription flotte sans point d'attache dans un fonds de seize
% mille feuillets.

\newcommand{\folder}[1]{\gdef\grothFolder{#1}}
\newcommand{\batch}[1]{\gdef\grothBatch{#1}}
\newcommand{\leaves}[2]{\gdef\grothFirst{#1}\gdef\grothLast{#2}}
\newcommand{\foldertitle}[1]{\gdef\grothFoldertitle{#1}}
\gdef\grothFolder{?}\gdef\grothBatch{?}\gdef\grothFirst{?}\gdef\grothLast{?}\gdef\grothFoldertitle{}

% --- Le résumé --------------------------------------------------------------
%
% Il ouvre la lecture modernisée, sous le titre « Résumé », et il est composé
% en petit. Ce n'est pas un ornement : c'est le seul passage du document qui
% ne soit pas des mathématiques, et un lecteur doit pouvoir le reconnaître
% avant de l'avoir lu — puis le sauter s'il connaît déjà le sujet, ou s'y
% arrêter s'il ne le connaît pas. Le filet qui le suit marque la reprise du
% corps.

\newenvironment{resume}
  {\small\setlength{\parskip}{0.45em}}
  {\par\medskip\hrule\medskip}

% --- Mots-clés ---------------------------------------------------------------
%
% En anglais, en clôture du résumé de la lecture modernisée : le vocabulaire
% moderne sous lequel chercher ce que les feuillets construisent. Le manifeste
% du site (`npm run manifest`) les extrait pour étiqueter la cote entière —
% cette ligne est la source unique des tags, il n'y a pas de fichier à côté.
\newcommand{\keywords}[1]{\par\smallskip\noindent
  {\footnotesize\textsc{Keywords} --- \textit{#1}}\par}

% --- L'appareil critique ----------------------------------------------------

% Le numéro de feuillet, en marge. C'est l'ancre qui permet de régler un
% désaccord sur une formule en regardant *un* feuillet plutôt qu'une chemise
% de six cents. Le site s'en sert aussi pour tourner les pages du fac-similé
% quand on fait défiler la transcription.
\newcommand{\leaf}[1]{%
  \marginnote{\footnotesize\color{gray!70}\textsf{#1}}%
  \label{leaf:#1}\ignorespaces}

% Illisible : on ne devine pas. Un mot inventé qui se lit comme les autres est
% le pire résultat possible.
\newcommand{\ill}{\textcolor{red!60!black}{[\,\ldots\,]}}

% Lecture douteuse : le mot est donné, et signalé comme incertain.
\newcommand{\uncertain}[1]{\uline{#1}}

% Ajout de l'éditeur — une lettre manquante, un mot sous-entendu.
\newcommand{\add}[1]{\textcolor{blue!50!black}{[#1]}}

% Biffé par Grothendieck lui-même. Ce qu'il a rayé fait partie du document :
% une rature dit souvent où la pensée a changé de direction.
\newcommand{\struck}[1]{\sout{#1}}

% Note du transcripteur, sur l'état matériel du feuillet.
\newcommand{\note}[1]{\par\smallskip{\small\color{gray!60!black}\textit{[#1]}}\par\smallskip}

% Note marginale de Grothendieck, distincte du corps du texte.
\newcommand{\marginal}[1]{\par\smallskip\noindent\hspace{1em}%
  {\small\color{gray!60!black}#1}\par\smallskip}

% --- Titre ------------------------------------------------------------------

\AtBeginDocument{%
  \begin{center}
    {\large\bfseries Fonds Alexandre Grothendieck --- cote n\textsuperscript{o}~\grothFolder}\\[2pt]
    {\itshape\grothFoldertitle}\\[4pt]
    {\small feuillets \grothFirst--\grothLast\ (lot \grothBatch)}\\[2pt]
    {\footnotesize\color{gray!60!black}Transcription établie d'après le fac-similé numérisé
     par l'Université de Montpellier.\\
     Les lectures incertaines sont soulignées, les illisibles notées [\,\ldots\,],
     les ajouts entre crochets.}
  \end{center}
  \vspace{1em}}

\endinput


\folder{158}
\batch{2}
\pages{21}{40}
\dating{[vers 1990-1991]}
\watermark{Édition de démonstration}
\foldertitle{[Autour des Dérivateurs] : notes manuscrites (s.d.)}

\begin{document}

\section*{Un monoïde dont le classifiant porte un objet 0-connexe non 1-connexe}

\note{le titre ci-dessus est de l'éditeur : ces feuillets poursuivent un raisonnement commencé aux feuillets précédents et ne portent aucun titre}

\note{les conditions sont numérotées sur les feuillets par des lettres et des chiffres entourés d'un cercle ; elles sont rendues ici (A), (B), \ldots{} et (0), (1), \ldots{}}

\page{21}\note{sa pagination : p. 3}

$\mathrm{Ob}\,X = F = M/R$.

Si $x, y \in \mathrm{Ob}\,X$, on a
\[
  \mathrm{Hom}(x,y) \simeq \{\, u \in M \mid x = yu \,\},
\]
la composition des flèches étant donnée par la composition de $M$. La catégorie
$X$ ayant $a = \dot{1}$ comme plus grand élément, \uncertain{dans notre sujet}
\ill{} de la forme
\[
  x = \dot{u} = au .
\]

\note{le point note le passage au quotient : $\dot{u}$ est la classe de $u$ dans $F = M/R$, et $a = \dot{1}$ ; la lecture est confirmée par le feuillet 31, où il écrit $e = \dot{1}$}

À quelles conditions $X$ \uncertain{admet-il} un plus petit élément $e$, lequel
sera de la forme $e = a\psi$ (pour $\psi \in M$ fixé) ?

\marginal{Notons que $X = Y_{/F}$}

\note{entre ces deux lignes court une rature où se lisent encore « pseudo-cofiltrant » et « il faut $F \in \mathrm{Ob}\,Y^\wedge$ »}

On voit tout de suite que la condition est la suivante :

(A) $\forall\, u \in M$, $\exists\, u' \in M$ tel que $uu' \overset{R}{\sim} \psi$,
(i.e. $\exists\, p,q \in \mathbb{N}$ \ill{} $\varphi^{p}\psi = \varphi^{q}uu'$).

Par exemple, il suffit que $\psi$ soit un élément minimal de $M$
(i.e. $\forall\, u$ $\exists\, u'$ \ill{} $\psi = uu'$).

\page{22}\note{sa pagination : p. 4}

Posons
\[
  M_0 = \mathrm{End}_X(e) = \{\, f \in M \mid ef = e,\ \text{i.e. } a\psi f = a\psi,
  \ \text{i.e. } \psi f \overset{R}{\sim} \psi \,\}.
\]

Ainsi, on a pour $f \in M$
\[
  f \in M_0 \iff \forall\, p,q \in \mathbb{N} \quad \varphi^{p}\psi f = \varphi^{q}\psi .
\]

\note{le quantificateur est bien tracé $\forall$, nettement distinct du $\exists$ en 3 renversé qu'il forme au bas du feuillet 21, alors que la suite (« ne dépend pas du choix de $p,q$ ») demande $\exists$}

Je vais écrire une condition sur $\varphi,\psi$ qui impliquera que pour
$f \in M_0$ fixé,
\[
  \delta(f) \overset{\text{déf}}{=} q - p
\]
ne dépend pas du choix de $p,q$.

\marginal{Notation}

Supposons donc qu'on ait
\[
  \begin{cases}
    \varphi^{p}\psi f = \varphi^{q}\psi \\
    \varphi^{p'}\psi f = \varphi^{q'}\psi
  \end{cases}
\]
on en conclut
\[
  \varphi^{p+p'}\psi f =
  \begin{cases}
    \varphi^{p'+q}\psi \\
    \varphi^{p+q'}\psi
  \end{cases}
\]
donc

\page{23}\note{sa pagination : p. 5}

\[
  \varphi^{p'+q}\psi = \varphi^{p+q'}\psi .
\]

\marginal{On veut en déduire $p'+q = p+q'$, i.e. $q-p = q'-p'$}

On va supposer

(B) Si $r,s \in \mathbb{N}$ sont tels que $\varphi^{r}\psi = \varphi^{s}\psi$, on a $r = s$.

On trouve alors une application
\[
  M_0 \xrightarrow{\ \delta\ } \mathbb{Z} ,
\]
et il est immédiat que celle-ci est \uncertain{multiplicative}, i.e.
$\delta(fg) = \delta(f) + \delta(g)$ :
\[
  \begin{cases}
    \varphi^{p}\psi f = \varphi^{q}\psi \\
    \varphi^{p'}\psi g = \varphi^{q'}\psi
  \end{cases}
\]
implique $\varphi^{p}\psi fg = \varphi^{q}\psi g$

\marginal{multiplier par $g$ à droite}

et en multipliant par $\varphi^{p'}$ à gauche
\[
  \varphi^{p+p'}\psi fg = \varphi^{q+p'}\psi g
  = \varphi^{q}\bigl(\underbrace{\varphi^{p'}\psi g}_{\varphi^{q'}\psi}\bigr)
  = \varphi^{q+q'}\psi
\]
donc $\delta(fg) = (q+q') - (p+p') = \delta(f) + \delta(g)$.

Donc $\delta$ se factorise en
\[
  M_0 \xrightarrow{\ \mathrm{can}\ } G_0 \xrightarrow{\ \delta_0\ } \mathbb{Z} .
\]

\note{sous $G_0$ il porte « groupe enveloppant de $M_0$ », et une flèche courbe étiquetée $\delta$ joint directement $M_0$ à $\mathbb{Z}$}

\page{24}\note{sa pagination : p. 6}

Je vais supposer

(C) L'homo\uncertain{morphisme} $M_0 \xrightarrow{\ \delta\ } \mathbb{Z}$ n'est
pas nul, i.e. $\exists\, f \in M_0$ et $p,q \in \mathbb{N}$, avec
$\varphi^{p}\psi f = \varphi^{q}\psi$ et $p \ne q$.

\struck{\ill{}}

\note{suit un bloc de six lignes, ouvert par « Exemple. Si $\psi \sim \ldots{} $ », entièrement biffé de longs traits horizontaux ; on y distingue encore $\psi f \sim \psi$, $\delta(f) = q-p$ et $\delta(\psi) = 1$, mais l'ensemble ne se laisse pas lire}

N.B. Dans le cas \ill{} objet \uncertain{minimal et maximal} \ill{}

La condition (C) implique que $G_0 = \pi_1(X_0,e)$ admet un quotient isomorphe à
$\mathbb{Z}$ (l'image de $G_0$ dans $\mathbb{Z}$ par $\delta_0$).

\marginal{$\varphi^{\alpha}\psi f = \varphi^{\beta}f$ \qquad $\psi f \overset{R}{\sim} \psi$ \qquad \ill{}}

\note{ces trois formules sont jetées au bas du feuillet, sans lien apparent avec ce qui précède ; la troisième ne se laisse pas lire}

\page{25}\note{sa pagination : p. 7}

\marginal{\ill{}}

\note{en tête du feuillet, un bloc de six lignes écrites en biais, dont on ne tire que « Si on \ldots{} », « N.B. pour \ldots{} » et « \ldots{} qu'il \ldots{} »}

\ill{} maintenant appliquer la prop. 1, pour calculer $\pi_1(X,e)$, en termes de
$\pi_1(X_0,e) = G_0$, et des relations explicitées par ailleurs. On veut
\struck{\ill{}} \add{montrer} que les images de ces relations dans $\mathbb{Z}$
par $\delta_0$ sont nulles, donc que $\pi_1(X_0,e)$ admet un groupe quotient
isomorphe à $\mathbb{Z}$, donc n'est pas $\simeq \{1\}$, donc que $X_0$ n'est pas
1-connexe.

\note{les deux indices $0$ sont bien tracés ici, alors que l'énoncé auquel le raisonnement aboutit, à sa p. 13 (feuillet 31), porte sur $\pi_1(X,e)$ et sur $X$}

Je reprends donc le diagramme ($\ast$)

\begin{tikzcd}
x \arrow[r, "f"] & y \arrow[r, "g"] & z \\
 & e \arrow[ul, "u"'] \arrow[u, "v"] \arrow[ur, "w"]
\end{tikzcd}

\note{sur le feuillet $e$ est placé sous $x$, les trois flèches partant de lui en éventail}

en posant
\[
  \begin{cases}
    z = a\rho \\
    y = a\rho g \\
    x = a\rho g f
  \end{cases}
  \qquad\text{donc}\qquad
  \begin{cases}
    e = xu = a\rho g f u \\
    e = yv = a\rho g v \\
    e = zw = a\rho w
  \end{cases}
\]
et comme $e = a\psi$, les trois relations précédentes s'écrivent
\[
  (\ast\ast) \qquad
  \boxed{\ \rho g f u \overset{R}{\sim} \rho g v \overset{R}{\sim} \rho w \overset{R}{\sim} \psi\ }
\]

La donnée d'un diagramme ($\ast$) dans $X$, avec $z = a\rho$ (pour $\rho$ fixé)
équivaut à celle des \uncertain{six} éléments

\page{26}\note{sa pagination : p. 8}

$\rho, f, g, u, v, w \in M$ de $M$, satisfaisant les relations ($\ast\ast$), ou
encore tels qu'il existe $p,q,r,s \in \mathbb{N}$ avec
\[
  \boxed{\ \varphi^{p}(\rho g f u) = \varphi^{q}(\rho g v) = \varphi^{r}(\rho w) = \varphi^{s}\psi\ }
\]

Ceci posé, on doit poser la relation
\[
  \boxed{\ [w,\, gv]\,[v,\, fu]\,[w,\, gfu]^{-1} = 1\ } ,
\]

Pour l'expliciter, considérons de façon générale

\begin{tikzcd}
e \arrow[r, "\lambda", bend left=25] \arrow[r, "\mu"', bend right=25] & \xi = a\sigma
\end{tikzcd}

donc $e = a\sigma\lambda = a\sigma\mu$, i.e.
\[
  \boxed{\ \sigma\lambda \overset{R}{\sim} \sigma\mu \overset{R}{\sim} \psi\ } ,
\]
et choisissons $\lambda', \mu' \in M$ tels que
$\lambda\lambda' = \mu\mu' \overset{\text{déf}}{=} \pi$, de sorte qu'on a un
diagramme commutatif

\begin{tikzcd}
 & e \arrow[dr, "\lambda"] \\
\eta \arrow[ur, "\lambda'"] \arrow[dr, "\mu'"'] & & \xi \\
 & e \arrow[ur, "\mu"']
\end{tikzcd}

donc
\[
  \eta =
  \begin{cases}
    e\lambda' = a\sigma\lambda\lambda' = a\sigma\pi \\
    e\mu' = a\sigma\mu\mu' = a\sigma\pi
  \end{cases}
\]

\note{sur le feuillet l'accolade se referme entre les signes « $=$ » et les produits $\lambda\lambda'$, $\mu\mu'$, qui y sont mis en facteur}

On va \struck{choisir} prendre une flèche $e \xrightarrow{\ \nu\ } \eta$
\ill{}, i.e. $\nu \in M$ tel que

\page{27}\note{sa pagination : p. 9}

$e = \eta\nu =$ \struck{\ill{}} $a\sigma\pi\nu$, i.e. tel que
\[
  \boxed{\ \sigma\pi\nu \overset{R}{\sim} \psi\ }
  \qquad\text{où } \pi = \lambda\lambda' = \mu\mu' .
\]

Alors on a
\[
  \begin{cases}
    \lambda'\nu,\ \mu'\nu \in M_0 \\
    \lambda(\lambda'\nu) = \mu(\mu'\nu)
  \end{cases}
\]

\note{entre $\lambda'\nu$ et $\mu'\nu$ une esperluette est biffée}

d'où
\[
  [\mu,\lambda] = (\mu'\nu)(\lambda'\nu)^{-1} .
\]

L'existence de $\nu$, satisfaisant $(\sigma\pi)\nu \overset{R}{\sim} \psi$, est
garantie par (A) : on obtient (pour tout $u \in M$) un $\theta(u)$ tel que
$u\,\theta(u) \overset{R}{\sim} \psi$, et on prendra ici
$\nu = \theta(\sigma\pi)$. Donc on a
\[
  [\mu,\lambda] = \bigl(\mu'\,\theta(\sigma\pi)\bigr)\bigl(\lambda'\,\theta(\sigma\pi)\bigr)^{-1}
\]
où $\lambda', \mu' \in M$ sont tels que
$\lambda\lambda' = \mu\mu'\ (\overset{\text{déf}}{=} \pi)$, et $\sigma$
\add{$\in M$} \struck{\ill{}} tel que $\xi = a\sigma$.

Je crois que je vais craquer, et restreindre la généralité, en supposant \ill{}
que $\psi$ est un objet minimal (ce qui implique (A)), de sorte que pour tout
\uncertain{$u$} il existe $\theta(u)$ tel que
\[
  u\,\theta(u) = \psi \qquad \text{(plus seulement } \overset{R}{\sim} \psi\text{)} ;
\]
dans ce cas \ill{} on peut prendre $\lambda' = \theta(\lambda)$,
$\mu' = \theta(\mu)$, et on aura
$\lambda\lambda' = \mu\mu' \overset{\text{déf}}{=} \pi = \psi$.

\page{28}\note{sa pagination : p. 10}

\[
  [\mu,\lambda] = \bigl(\underbrace{\theta(\mu)\,\theta(\sigma\psi)}_{\in M_0}\bigr)
  \bigl(\underbrace{\theta(\lambda)\,\theta(\sigma\psi)}_{\in M_0}\bigr)^{-1}
\]

On \ill{} aussi \ill{} les quantités analogues qui figurent dans les relations de
la \ill{}

\note{la ligne porte ensuite deux signes que l'on peut lire $\forall\,\exists$, puis « 8 : » — vraisemblablement un renvoi à sa page 8, le feuillet 26, où la relation est posée}

\[
  \begin{array}{l}
    [w,\, gv] = \bigl(\theta(w)\,\theta(\rho\psi)\bigr)\bigl(\theta(gv)\,\theta(\rho\psi)\bigr)^{-1} \\
    {}[v,\, fu] = \bigl(\theta(v)\,\theta(\rho g\psi)\bigr)\bigl(\theta(fu)\,\theta(\rho g\psi)\bigr)^{-1} \\
    {}[w,\, gfu] = \bigl(\theta(w)\,\theta(\rho\psi)\bigr)\bigl(\theta(gfu)\,\theta(\rho\psi)\bigr)^{-1}
  \end{array}
\]

La relation devient dès lors
\[
  \boxed{\ \bigl(\theta(gv)\,\theta(\rho\psi)\bigr)^{-1}
  \bigl(\theta(v)\,\theta(\rho g\psi)\bigr)
  \bigl(\theta(fu)\,\theta(\rho g\psi)\bigr)^{-1}
  \bigl(\theta(gfu)^{-1}\,\theta(\rho\psi)\bigr) = 1\ }
\]

\note{le dernier facteur est renvoyé à la ligne du dessous : le « $-1$ » y est porté juste après la parenthèse qui ferme $\theta(gfu)$, et la parenthèse ouvrante du facteur n'est pas refermée sur le feuillet. La règle de composition de ses p. 9-10 donnerait ici $\bigl(\theta(gfu)\,\theta(\rho\psi)\bigr)$, sans exposant}

( N.B. $w$ y a disparu sans laisser de traces. ) En fait, on a besoin de s'en
\ill{} un peu plus, pour \ill{} ; \ill{}

\page{29}\note{sa pagination : p. 11}

On va supposer

\marginal{ce qui implique}

(D) $\boxed{\ \psi \overset{R}{\sim} 1\ }$ i.e. $\exists\, \alpha,\beta \in \mathbb{N}$
\ill{} $\boxed{\ \varphi^{\alpha}\psi = \varphi^{\beta}\ }$

\note{la lettre entourée se lit mal ; la suite (A), (B), (C) des conditions donne (D)}

Dans les relations
\[
  \rho g f u \overset{R}{\sim} \rho g v \overset{R}{\sim} \rho w \overset{R}{\sim} \psi
\]
équivalent :
\[
  \rho g f u \overset{R}{\sim} \rho g v \overset{R}{\sim} \rho w \overset{R}{\sim} 1
\]
i.e. :
\[
  \boxed{\ \rho g f u,\ \rho g v,\ \rho w \in M_0\ }
\]

(Que $w$ ait disparu de la relation fondamentale ne doit pas nous inquiéter,
c'est une relation en \struck{$f, \psi, g, q$} $\rho, f, g, u, v$, et on peut
poser $w = gv$, de sorte qu'il ne reste, comme conditions sur ces cinq variables
ou quantités, que
\[
  \underbrace{\rho g f u}_{\overset{\text{déf}}{=}\ \pi},\qquad
  \underbrace{\rho g v}_{\overset{\text{déf}}{=}\ \pi'} \ \in M_0 ,
\]
donc on trouve deux éléments $\pi, \pi' \in M_0$.

\note{la parenthèse ouverte à « (Que $w$ ait disparu » n'est pas refermée. $\pi$ est ici redéfini : au feuillet 26 il désignait $\lambda\lambda' = \mu\mu'$}

Ceci posé, on trouve les relations

\page{30}\note{sa pagination : p. 12}

\[
  \pi\bigl(\underbrace{\theta(gv)\,\theta(\rho\psi)}_{c_1}\bigr)
  = \pi\bigl(\underbrace{\theta(v)\,\theta(\rho g\psi)}_{c_2}\bigr) = \psi
\]
\[
  \pi'\bigl(\underbrace{\theta(fu)\,\theta(\rho g\psi)}_{c_3}\bigr)
  = \pi'\bigl(\underbrace{\theta(gfu)\,\theta(\rho\psi)}_{c_4}\bigr) = \psi
\]

\note{les deux lignes sont réunies par une accolade sur le feuillet, et sous le premier $\pi$ il rappelle, en le soulignant d'un double trait, $\pi = \rho g f u$}

d'où
\[
  \begin{array}{l}
    \delta(c_1) = \delta(c_2) = \delta(\psi) - \delta(\pi) \\
    \delta(c_3) = \delta(c_4) = \delta(\psi) - \delta(\pi')
  \end{array}
\]
d'où
\[
  \delta_0\bigl(c_1 c_2^{-1} c_3 c_4^{-1}\bigr)
  = \underbrace{\bigl(\delta(c_1) - \delta(c_2)\bigr)}_{0}
  + \underbrace{\bigl(\delta(c_3) - \delta(c_4)\bigr)}_{0} = 0
\]
donc \struck{$G_0 \xrightarrow{\ \delta\ }$} $\delta_0 : G_0 \longrightarrow \mathbb{Z}$
se factorise en un homomorphisme surjectif
\[
  \delta'_0 : \pi_1(X,e) \longrightarrow \mathbb{Z} .
\]

\page{31}\note{sa pagination : p. 13}

En résumé :

\textbf{Proposition 2.} Soient $M$ un monoïde, \struck{\ill{}}
\struck{et supposons} $\varphi$ et $\psi$ deux éléments de $M$. On suppose :

a) $\psi$ est un élément minimal de $M$, i.e. pour tout $u \in M$, existe
$\theta(u) \in M$ avec $\psi = u\,\theta(u)$.
\add{(Donc $M$, i.e. $Y = B_M$, est pseudo-cofiltrant, i.e. $Y^\wedge$ est
loc. irréductible.)}

b) $\psi \overset{R_\varphi}{\sim} 1$ pour la relation d'équivalence \ill{}
$R_\varphi$ définie par $\varphi$
\[
  \bigl(u \overset{R_\varphi}{\sim} v \iff \exists\, p,q \in \mathbb{N}
  \ \ill{}\ \varphi^{p}u = \varphi^{q}v\bigr),
\]
i.e. $\exists\, \alpha,\beta \in \mathbb{N}$ \ill{}
$\varphi^{\alpha}\psi = \varphi^{\beta}$\struck{$\psi$}.

c) Si $r,s \in \mathbb{N}$ sont tels que $\varphi^{r}\psi = \varphi^{s}\psi$,
alors $r = s$.

Soit $Y = B_M$, $F = M/R_\varphi$ considéré comme $M$-ensemble, donc i.e. comme
objet de $Y^\wedge$, $X = Y_{/F}$. Alors l'objet $e = \dot{1}$ de $X$ est à la
fois le plus \struck{petit} \add{grand} et un plus petit élément de $X$. De plus,
\struck{$\pi_1(X,e)$} il existe un homomorphisme surjectif
\[
  \pi_1(X,e) \longrightarrow \mathbb{Z} ,
\]
en particulier $X$ n'est pas 1-connexe.

\page{32}\note{sa pagination : p. 14}

\textbf{Corollaire.} Supposons que $M$ soit \struck{\ill{}} isomorphe, soit au
monoïde $\mathrm{Ep}(E)$ où $E$ est un ens. infini, soit au monoïde opposé :
$\mathrm{Mon}(E)$, où $E$ un ens. \ill{} infini.

\note{ni $\mathrm{Ep}(E)$ ni $\mathrm{Mon}(E)$ ne sont définis ici ; le feuillet 36 emploie $\mathrm{Ep}(M_0)$ pour un monoïde d'applications $M_0 \to M_0$}

Alors $M$ est pseudo-cofiltrant \struck{\ill{}} (il admet \struck{\ill{}} des
éléments minimaux, mais les \ill{} vérifiant la loi \uncertain{$fuf = uf$}
$\Longrightarrow u = v$), $Y = B_M$ est 1-connexe, mais il existe
$F \in \mathrm{Ob}\,Y^\wedge$ tel que $X = Y_{/F}$ soit 0-connexe et non
1-connexe.

Pour ceci, il suffit de prendre \struck{pour} la construction de la proposition,
en prenant \struck{pour} $\psi = \varphi$ n'importe quel \uncertain{objet}
minimal de $M$. Il suffit de vérifier qu'on ne peut avoir
$\varphi^{r} = \varphi^{s}$ avec $r \ne s$. Mais si $r = s+h$, $h > 0$, on aurait
$\varphi^{h}\varphi^{s} = \mathrm{id}\cdot\varphi^{s}$, donc $\varphi$ serait
inversible.

\note{le premier exposant peut aussi se lire $p$ ; c'est la condition « avec $r \ne s$ » qui a fixé la lecture}

N.B. $\varphi^{h} = 1$,

\page{33}\note{sa pagination : p. 15}

Mais un élément inversible ne peut être minimal. En effet, on aurait pour tout
$u \in M$ un $u' \in M$ tel que $\varphi = uu'$, \struck{\ill{}} et en posant
$v = u'\varphi^{-1}$, on aurait
\[
  uv = 1 ;
\]
montrons que l'on aurait aussi
\[
  p \overset{\text{déf}}{=} vu = 1 ,
\]
en effet on a $p^{2} = p$, i.e. $p\cdot p = \mathrm{id}\cdot p$ dans \ill{} $M$,
$p = \mathrm{id}$. Donc tous les objets de $M$ seraient inversibles, i.e. $M$
serait un groupe, ce qui n'est pas le cas.

N.B. Considérons \struck{\ill{}} \struck{une partie} (sous-$M$-ensemble :
\struck{\ill{}} \uncertain{idéaux à} droite) $F$ de $M$, donc $F$ \ill{} regardés
comme des objets de $Y^\wedge$ ($Y = B_M$), d'où une \ill{} $X = Y_{/F}$. On peut
montrer \ill{} et pourtant les $X$ sont tous 1-connexes, \ill{} que j'ai un peu
\ill{} mal à construire un $F$ \struck{\ill{}} \add{0-connexe} qui ne soit
1-connexe !

\note{les feuillets 35 et 36 sont reliés dans l'ordre inverse du sien : le feuillet 36 porte « 15 bis » et s'arrête au milieu d'une phrase que le feuillet 35, numéroté « 16 », reprend. Le raisonnement se lit donc 33, 36, 35, 37, 39}

\page{35}\note{sa pagination : p. 16 ; ce feuillet reprend la phrase interrompue au bas du feuillet 36, numéroté « 15 bis »}

$\overline{\varphi}^{\,h} = 1$, donc $\overline{\varphi}$ est inversible. Mais
alors \struck{\ill{}}
$\overline{\psi} \overset{\text{déf}}{=} \overline{\varphi}^{\,\beta}$ est
inversible (cet \ill{} d'ordre fini), \struck{\ill{}} et d'autre part il est
minimal comme \ill{} image de $\overline{\varphi}^{\,\beta}$ qui est minimal.
\struck{Ainsi}

Mais si dans un monoïde $\overline{M}$ \struck{\ill{}} \ill{} on a la loi
\[
  \bigl[\ \bar{u}\bar{f} = \bar{v}\bar{f} \Longrightarrow \bar{u} = \bar{v}\ \bigr],
\]
on a \ill{} un élément minimal et inversible, \ill{} \ill{} que $\overline{M}$
est un groupe. Car $\forall\, u \in \overline{M}$,
$\exists\, v \in \overline{M}$ \ill{} tels que $\sigma = uv$, i.e.
$u(v\sigma^{-1}) = 1$, et alors $p = vu$ \uncertain{satisfait} $p^{2} = p$
\struck{donc} $p\cdot p = 1\cdot p$, i.e. $p = 1$, donc $u$ \ill{} \ill{} est
inversible.

\note{dans « $p \cdot p = 1 \cdot p$ », le $1$ est écrit par-dessus un $\mathrm{id}$ biffé}

On a prouvé :

\textbf{Proposition 3.} Soit $M$ un monoïde ayant un \struck{objet} élément
minimal $\psi$, et soit $\varphi \in M$ tel qu'il existe $\beta \in \mathbb{N}$
avec $\varphi^{\beta}$ minimal (et \ill{} pas exact que $\varphi$ minimal
\struck{\ill{}}, par ex. \ill{} $\varphi = \psi$). S'il existe
$r,s \in \mathbb{N}$ avec $r \ne s$ et $\varphi^{r}\psi = \varphi^{s}\psi$,

\note{l'énoncé s'arrête là : le bas du feuillet est occupé par une grande boucle et la conclusion manque}

\page{36}\note{sa pagination : « 15 bis » ; ce feuillet précède, dans son ordre, celui qui porte le numéro 16 (feuillet 35)}

\struck{\ill{}}

\note{le feuillet s'ouvre sur une ligne entièrement biffée, où l'on reconnaît la relation du feuillet 28 recopiée avec $\alpha$ à la place de $\theta$ : $(\alpha w)^{-1}(\alpha gv)(\alpha gv)^{-1}(\alpha gfu)\bigl((\alpha w)^{-1}(\alpha gfu)\bigr)^{-1} = 1$}

(0) $M$ pseudo-cofiltr. \ill{}, i.e. $\forall\, u,v$ $\exists\, u',v'$ \ill{}
$uu' = vv'$

(1) $\forall\, u$, $\exists\, u'$ tel que $uu' \overset{R}{\sim} \psi$, i.e.
$\exists\, p,q \in \mathbb{N}$ tels que $\varphi^{p}uu' = \varphi^{q}\psi$

(2) $\exists\ a \xrightarrow{\ \lambda\ } e = a\psi$, i.e. $\exists\, \lambda$
tel que $\psi\lambda \overset{R}{\sim} 1$, i.e.
$\exists\, \alpha,\beta \in \mathbb{N}$ tels que
$\boxed{\ \varphi^{\alpha}\psi\lambda = \varphi^{\beta}\ }$

[ N.B. Si $\psi$ est \emph{minimal}, alors (0) et (1) sont satisfaits, et (2) est
satisfait ssi \struck{\ill{}} $\exists\, \beta \in \mathbb{N}$ tel que
$\varphi^{\beta}$ soit minimal (. Il \ill{} pour un tel $\beta$,
$\exists\, \lambda \in M$ \ill{}) pour un tel $\beta$, tel que
$\psi\lambda = \varphi^{\beta}$ ]

(3) $\forall\, r,s \in \mathbb{N}$ avec $\varphi^{r}\psi = \varphi^{s}\psi$, on a
$r = s$.

[ N.B. Si $\psi$ est minimal, alors la relation
\[
  R_0 : \qquad u \overset{R_0}{\sim} v \iff u\psi = v\psi
\]
est aussi la relation d'équivalence définie par l'application
\[
  \begin{array}{rcl}
    M &\longrightarrow& \mathrm{Ep}(M_0) \\
    u &\longmapsto& u_{M_0} = \bigl(\lambda \mapsto u\lambda : M_0 \to M_0\bigr)
  \end{array}
\]
(puisque $u\psi = v\psi \Longrightarrow u\psi\alpha = v\psi\alpha$
$\forall\, \alpha \in M$, or $\psi M = M_0$, i.e. l'ens. des $\psi\alpha$ est
l'ens. des $\lambda \in M_0$). Soit
$\overline{M} \simeq M/R_0 \subset \mathrm{Ep}(M_0)$ le monoïde image.
\struck{Ainsi} C'est un monoïde \ill{} dans lequel on a :
$\bar{u}\bar{f} = \bar{v}\bar{f} \Longrightarrow \bar{u} = \bar{v}$, et il est
pseudo-cofiltrant, comme image de $M$. ($X$ n'est filtrant que s'il est réduit à
1.) Ainsi la relation $\varphi^{r}\psi = \varphi^{s}\psi$ dans (3) équivaut à
$\overline{\varphi}^{\,r} = \overline{\varphi}^{\,s}$, et $\overline{\varphi}$
est l'image de $\varphi$ dans $\overline{M}$. Si on a $r = s+h$, $h \ge 1$,
\struck{\ill{}} :

\page{37}\note{sa pagination : p. 17}

donc l'image $\overline{M}$ de $M$ dans $\mathrm{Ep}(E)$ ($E$ = ens. des éléments
minimaux de $M$) est un groupe. Donc si $Y = B_M$ est 1-connexe, ce groupe est
réduit à $\{1\}$, ce qui signifie que pour \struck{\ill{}} $\lambda \in E$
(p. ex. pour $\lambda = \psi$), \struck{\ill{}} et $u \in M$, i.e. \ill{}
$u\lambda = \lambda$, ce qui implique que $M$ est cofiltrant.

\note{la lettre $E$ est tracée d'un double trait sur le feuillet}

Dans le cas contraire
($\varphi^{r}\psi = \varphi^{s}\psi \Longrightarrow r = s$) \add{\ill{} immédiat
\ill{} dans $\overline{M}$}, prenant la relation d'équivalence $R_\varphi$
\struck{définie} par
\[
  u \overset{R_\varphi}{\sim} v \iff \exists\, p,q \in \mathbb{N}
  \ \text{avec}\ \varphi^{p}u = \varphi^{q}v ,
\]
et \ill{} le quotient $F = M/R_\varphi$, considéré comme $M$-ens. à droite
\add{(pseudo-cofiltrant)} i.e. comme objet de $Y^\wedge$, alors $Y_{/F}$ est
0-connexe (tous ses objets sont \uncertain{équivalents}), mais elle n'est pas
1-connexe, plus précisément son groupe fondamental est isomorphe à
$\pi_1(M_0)$, où
\[
  M_0 = \{\, f \in M \mid \psi f \overset{R_\varphi}{\sim} \psi \,\} ,
\]
et admet un groupe quotient isomorphe à $\mathbb{Z}$, à savoir : l'homomorphisme
surjectif
\[
  \pi_1(M_0) = G_0 \xrightarrow{\ \delta_0\ } \mathbb{Z}
\]
\ill{} défini sur $M_0$ par
\[
  \delta(f) = q - p \qquad \text{si } \varphi^{p}\psi f = \varphi^{q}\psi .
\]

[ N.B. On \struck{\ill{}} \ill{} comme $\psi \in M_0$, on a $\varphi\psi \in M_0$,
\struck{\ill{}} donc $\exists\, f$ tel que $\psi f = \varphi\psi$, et donc
$f \in M_0$ et $\delta(f) = 1$ ] Donc on a $H^{1}(X,\mathbb{Z}) \ne 0$.

\page{39}\note{sa pagination : p. 18}

\struck{Donc} \add{Ainsi}

\textbf{Corollaire.} Soit $M$ un monoïde ayant un objet minimal. Si
\struck{$M$ n'est pas} \ill{} $Y = B_M$ est tel que tous les objets 0-connexes de
$Y^\wedge$ sont 1-connexes (ou simplement, que leur $H^{1}(X,\mathbb{Q}) = 0$),
donc $M$ est cofiltrant.

\textbf{Théorème.} Soit $X$ \struck{une catégorie} \add{$\in \mathrm{Cat}$}
pseudo-cofiltrante \ill{} vérifiant : les objets 0-connexes de $X^\wedge$
(hyp. vérifiée si les objets 0-connexes de \ill{} sont 1-connexes, ou \ill{}
seulement il existe un groupe \add{abélien} $h \ne 0$ tel que pour \ill{}
$H^{1}(F,h) \ne 0$).

On suppose que $X$ admet \struck{des objets} un plus \struck{grand} petit objet
$e$, et que $X_{/e}$ (qui est \struck{\ill{}} pseudo-cofiltrant) en admet un
\uncertain{également}, \ill{} i.e. ce qui revient au même, que le monoïde
$M = \mathrm{End}_X(e)$ admet un \struck{objet} élément minimal $\varphi$.
(N.B. cette hypothèse de minimum est satisfaite si \uncertain{$\mathrm{Ob}\,X$}
et $M = \mathrm{End}_X(e)$ sont finis, en particulier si $X$ est
\uncertain{fini}.)

\note{dans « $M = \mathrm{End}_X(e)$ », un premier « End » est biffé et récrit}

Considérons l'opération de $\mathbb{N}$ sur $e$ (identifié à un objet de
$X^\wedge$) définie par $\varphi$. \uncertain{$\mathbb{N}$ opère par $\varphi$},
et distinguons deux cas :

a) Cette opération \struck{se} fait « sans pts fixes », i.e.
\[
  \mathbb{N} \times e \longrightarrow e \times e ,\qquad
  (n,u) \longmapsto (\varphi^{n}u,\, u)
\]
est mono. Il revient au \ill{} de dire que pour

\end{document}
